% GENERATED FILE. Edit canonical lessons, then run npm run generate:books.
% canonical-source-hash: fnv1a64:952cb671c2667dc4
% canonical-lessons: HI-C230-jhukna, HI-C230-latakna, HI-C230-latkana, HI-C230-tairana, HI-C230-dubna

\chapter{To bend down, To hang, To hang, To float, To sink}
\label{ch:hi-to-bend-down-to-hang-to-hang-to-float-to-sink}
\hlchaptermodality{230}

\begin{chapteropening}
\textbf{By the end of this chapter:} \emph{I can say to bend down, to hang, to hang, to float and to sink.}

Complete the last lesson of chapter 230: I can say to bend down, to hang, to hang, to float and to sink.
\end{chapteropening}

\section[\texorpdfstring{\hi{झुकना} (jhuknā)}{झुकना (jhuknā)}]{\hi{झुकना} (jhuknā) — to bend down, to bow}

\label{lesson:HI-C230-jhukna}



\begin{quote}
Before the new one: say the Hindi for to spread, then the Hindi for to be left behind.
\end{quote}

\subsection*{You'll want to know: \hi{झुकना}}
\textbf{\hi{झुकना}} — \emph{jhuknā} — \textquotedblleft{}to bend down, to bow\textquotedblright{}.

\begin{cousinweb}
\textbf{\hi{झुकना}} (\emph{jhuknā}) is a tadbhava, a word inherited and reshaped by sound change. It comes from \emph{*jhukkati}, \textquotedblleft{}stoops, breaks\textquotedblright{}, a form CDIAL reconstructs from its descendants (CDIAL 5399).
\end{cousinweb}

\begin{grammarlens}[title={What you've built}]
One more word for this chapter.
\end{grammarlens}

\subsection*{Your turn}
\begin{itemize}
\raggedright
  \item \emph{Say it:} \emph{jhuknā}
  \item \emph{Say it:} \emph{jhuknā}, once more
  \item \emph{Say it:} say \emph{jhuknā}
  \item \emph{Recall:} say the Hindi for to spread, then the Hindi for to be left behind, then say \emph{jhuknā} again
\end{itemize}

\begin{tcolorbox}[breakable,colback=teal!4,colframe=teal!35!black,title={Before you move on}]
What does \hi{झुकना} mean? (\textquotedblleft{}To bend down, to bow\textquotedblright{}.) Say it once more.
\end{tcolorbox}
\section[\texorpdfstring{\hi{लटकना} (laṭaknā)}{लटकना (laṭaknā)}]{\hi{लटकना} (laṭaknā) — to hang}

\label{lesson:HI-C230-latakna}



\begin{quote}
Before the new one: say the Hindi for to be left behind, then the Hindi for to bend down.
\end{quote}

\subsection*{You'll want to know: \hi{लटकना}}
\textbf{\hi{लटकना}} — \emph{laṭaknā} — \textquotedblleft{}to hang\textquotedblright{}.

\begin{cousinweb}
\textbf{\hi{लटकना}} (\emph{laṭaknā}) probably belongs with a reconstructed \emph{*laṭṭa}, \textquotedblleft{}hanging, tangled\textquotedblright{} (CDIAL 10918). CDIAL lists Sindhi \emph{laṭkaṇu}, \textquotedblleft{}to hang, dangle\textquotedblright{}, and Gujarati \emph{laṭakvũ} there, but not the Hindi verb.
\end{cousinweb}

\begin{grammarlens}[title={What you've built}]
One more word for this chapter.
\end{grammarlens}

\subsection*{Your turn}
\begin{itemize}
\raggedright
  \item \emph{Say it:} \emph{laṭaknā}
  \item \emph{Say it:} \emph{laṭaknā}, once more
  \item \emph{Say it:} say \emph{laṭaknā}
  \item \emph{Recall:} say the Hindi for to be left behind, then the Hindi for to bend down, then say \emph{laṭaknā} again
\end{itemize}

\begin{tcolorbox}[breakable,colback=teal!4,colframe=teal!35!black,title={Before you move on}]
What does \hi{लटकना} mean? (\textquotedblleft{}To hang\textquotedblright{}.) Say it once more.
\end{tcolorbox}
\section[\texorpdfstring{\hi{लटकाना} (laṭkānā)}{लटकाना (laṭkānā)}]{\hi{लटकाना} (laṭkānā) — to hang (something)}

\label{lesson:HI-C230-latkana}



\begin{quote}
Before the new one: say the Hindi for to bend down, then the Hindi for to hang.
\end{quote}

\subsection*{You'll want to know: \hi{लटकाना}}
\textbf{\hi{लटकाना}} — \emph{laṭkānā} — \textquotedblleft{}to hang (something)\textquotedblright{}.

\begin{cousinweb}
\textbf{\hi{लटकाना}} (\emph{laṭkānā}) is the causing form of \emph{laṭaknā}, \textquotedblleft{}to hang\textquotedblright{}. It probably belongs with a reconstructed \emph{*laṭṭa}, \textquotedblleft{}hanging, tangled\textquotedblright{}, where CDIAL lists Sindhi \emph{laṭkaṇu}, \textquotedblleft{}to dangle\textquotedblright{} (CDIAL 10918), but not the Hindi verbs.
\end{cousinweb}

\begin{grammarlens}[title={What you've built}]
One more word for this chapter.
\end{grammarlens}

\subsection*{Your turn}
\begin{itemize}
\raggedright
  \item \emph{Say it:} \emph{laṭkānā}
  \item \emph{Say it:} \emph{laṭkānā}, once more
  \item \emph{Say it:} say \emph{laṭkānā}
  \item \emph{Recall:} say the Hindi for to bend down, then the Hindi for to hang, then say \emph{laṭkānā} again
\end{itemize}

\begin{tcolorbox}[breakable,colback=teal!4,colframe=teal!35!black,title={Before you move on}]
What does \hi{लटकाना} mean? (\textquotedblleft{}To hang (something)\textquotedblright{}.) Say it once more.
\end{tcolorbox}
\section[\texorpdfstring{\hi{तैराना} (tairānā)}{तैराना (tairānā)}]{\hi{तैराना} (tairānā) — to float (something)}

\label{lesson:HI-C230-tairana}



\begin{quote}
Before the new one: say the Hindi for to hang, then the Hindi for to hang.
\end{quote}

\subsection*{You'll want to know: \hi{तैराना}}
\textbf{\hi{तैराना}} — \emph{tairānā} — \textquotedblleft{}to float (something)\textquotedblright{}.

\begin{cousinweb}
\textbf{\hi{तैराना}} (\emph{tairānā}) is the causing form of \emph{tairnā}, \textquotedblleft{}to swim, float\textquotedblright{}, which CDIAL lists under Sanskrit \emph{tirate}, \textquotedblleft{}crosses over, sails across\textquotedblright{} (CDIAL 5821).
\end{cousinweb}

\begin{grammarlens}[title={What you've built}]
One more word for this chapter.
\end{grammarlens}

\subsection*{Your turn}
\begin{itemize}
\raggedright
  \item \emph{Say it:} \emph{tairānā}
  \item \emph{Say it:} \emph{tairānā}, once more
  \item \emph{Say it:} say \emph{tairānā}
  \item \emph{Recall:} say the Hindi for to hang, then the Hindi for to hang, then say \emph{tairānā} again
\end{itemize}

\begin{tcolorbox}[breakable,colback=teal!4,colframe=teal!35!black,title={Before you move on}]
What does \hi{तैराना} mean? (\textquotedblleft{}To float (something)\textquotedblright{}.) Say it once more.
\end{tcolorbox}
\section[\texorpdfstring{\hi{डूबना} (ḍūbnā)}{डूबना (ḍūbnā)}]{\hi{डूबना} (ḍūbnā) — to sink, to drown}

\label{lesson:HI-C230-dubna}



\begin{quote}
Before the new one: say the Hindi for to hang, then the Hindi for to float.
\end{quote}

\subsection*{You'll want to know: \hi{डूबना}}
\textbf{\hi{डूबना}} — \emph{ḍūbnā} — \textquotedblleft{}to sink, to drown\textquotedblright{}.

\begin{cousinweb}
\textbf{\hi{डूबना}} (\emph{ḍūbnā}) is a tadbhava, a word inherited and reshaped by sound change. It comes from \emph{*ḍubb}, \textquotedblleft{}sink\textquotedblright{}, a form CDIAL reconstructs from its descendants (CDIAL 5561). CDIAL explains it as a Middle Indo-Aryan \emph{buḍḍaï} with its consonants swapped.
\end{cousinweb}

\begin{grammarlens}[title={What you've built}]
One more word for this chapter.
\end{grammarlens}

\subsection*{Your turn}
\begin{itemize}
\raggedright
  \item \emph{Say it:} \emph{ḍūbnā}
  \item \emph{Say it:} \emph{ḍūbnā}, once more
  \item \emph{Say it:} say \emph{ḍūbnā}
  \item \emph{Recall:} say the Hindi for to hang, then the Hindi for to float, then say \emph{ḍūbnā} again
\end{itemize}

\begin{tcolorbox}[breakable,colback=teal!4,colframe=teal!35!black,title={Before you move on}]
What does \hi{डूबना} mean? (\textquotedblleft{}To sink, to drown\textquotedblright{}.) Say it once more.
\end{tcolorbox}
